\documentclass[10pt,a4paper]{amsart}
\usepackage[margin=19mm]{geometry}
\usepackage{amsmath,amssymb,mathtools}
\usepackage[scr=rsfs,cal=euler]{mathalfa}
\usepackage{xfrac}
\usepackage[unicode,hidelinks,bookmarks=false]{hyperref}
\newcommand{\ie}{i.e.\ }
\newcommand{\cf}{cf.\ }
\newcommand{\resp}{respectively\ }
\newcommand{\Subsection}{\subsection}
\providecommand{\nobreakdash}{\nobreak\hbox{-}}
\setlength{\emergencystretch}{2em}
\multlinegap0pt
\usepackage{amssymb}
\usepackage{tikz}
\usepackage{mathtools}
\newtheorem{lemma}{Lemma}
\renewcommand{\Im}{{\rm Im}\,}
\newcommand{\R}{\mathcal{R}}
\renewcommand{\Re}{{\rm Re}\,}
\renewcommand{\S}{\mathcal{S}}
\newcommand{\U}{\mathcal{U}}
\newcommand{\cd}{\cdot}
\newcommand{\cp}{\mathbb{C}} %<--- because of my last name
\newcommand{\ep}{\varepsilon}
\newcommand{\ft}{\widehat}
\newcommand{\ga}{\gamma}
\newcommand{\la}{\lambda}
\newcommand{\mt}{\mapsto}
\newcommand{\ord}{{\rm ord}}
\newcommand{\ov}[1]{\overline{#1}}
\newcommand{\p}{\varphi}
\renewcommand{\r}{\mathbb{R}}
\newcommand{\ra}{\rightarrow}
\DeclareMathOperator{\supp}{\mathrm{supp}}
\newcommand{\w}{\ov{w}}
\title{Classification of Fourier summation formulas on a horizontal strip}
\author{Guilherme Vedana}
\date{Source: arXiv:2608.10121v2 (CC BY 4.0)}
\begin{document}
\maketitle

\noindent\emph{Source and licence.} Classification of Fourier summation formulas on a horizontal strip. Version arXiv:2608.10121v2 (CC BY 4.0), distributed by arXiv under the Creative Commons Attribution 4.0 International licence, http://creativecommons.org/licenses/by/4.0/, as recorded in the arXiv licence metadata for this exact version and shown on the abstract page https://arxiv.org/abs/2608.10121v2. Full licence notice: \texttt{licenses/arxiv-2608.10121-CC-BY-4.0.txt}.\par\smallskip

\noindent\textit{Editorial scope.} Counted unit: the article's lemma thm:bridge (source0.tex lines 917--926). The article's complete original proof of it is delivered with it (source0.tex lines 927--975, one \texttt{proof} environment of 49 lines). No mathematics is written, repaired or re-derived: the statement and the proof are the article's own lines, in the article's own notation, and no retained line is substituted or re-flowed.  Retained uncounted as the source-local context the proof needs: lines 868--879.  Nothing was omitted from any retained span, so the package carries no omission record.  No retained line contains a citation, so the wrapper carries no bibliography and no bibliography entry.  No retained line contains a figure environment, an image-inclusion command or drawing code, so no image is delivered and the package declares no asset dependency.  Editorial formatting only, in addition: an AMS \texttt{amsart} preamble that restates the article's own declarations and macros used by the retained lines, namely the environments \texttt{lemma} on separate counters, together with the standard packages those lines need.  The retained environments print in this document as Lemma 1 in order of appearance, whereas the article prints the same items with the numbering of its own full document.  The complete native source of the article as distributed in the arXiv e-print for this exact version is bundled unchanged as \texttt{sources/207302/source0.tex}.
\begin{lemma}
    \label{lemma:asymptotic_decay}
    The following statements hold:
    \begin{itemize}
        \item[i)] Let $f_1,f_2\in L^1(\r)$ such that $|f_1(s)|\leq\frac{C}{(1+s^2)^n}$ and $|f_2(s)|\leq\frac{C}{(1+s^2)^n}$ whenever $|s|>R_1$, for some integer $n\geq1$. Then $|f_1*f_2(s)|\leq\frac{C_1}{(1+s^2)^n}$ for all $|s|\geq2R_1$.
        \item[ii)] Consider the region $\R:=\{z\in\cp^+;|\Re(z)|<1 \text{ and } \Im(z)>1\}$. Then, for any $\zeta\in\U$, $z,w\in\R$, and $T>1$, it holds that
        \begin{equation}
        \label{eq:unif_estimate_lemma}
        \left|\ft{G_k}(w,z,\cd)* \S_T(\zeta)\right|\ll\frac{1}{(1+(\Re{\zeta})^2)^{k+1}}.
        \end{equation}
    \end{itemize}
\end{lemma}

\begin{lemma}[The Bridge Lemma]
\label{thm:bridge}
If $(\mu=\nu+\eta,a)$ is an FS-pair on a strip such that $\ord(\mu)\leq 2(k+1)$ and $A\subset \{z;|\Im{z}|<1/2\}$, then
\begin{align}
\label{eq:bridge}
\lim_{T\ra\infty} \sum_{|\la|\leq T(k+1)} a(\la)G_k(w,z,\la){\ft{S_k}\left(\frac{\la}{T}\right)}&=\frac{1}{2\pi^{k+1}i}\int_{\r} \frac{1}{(t-z)(t-\w)}\cd\frac{d\nu(t)}{(1+t^2)^k}\\
&+\frac{1}{2\pi^{k+1}i}\sum_{\ga\in A} \frac{1}{(\ga-z)(\ga-\w)}\cd\frac{b(\ga)}{(1+\ga^2)^k}
\end{align}
for $z,w \in \cp^+$ in the region $\R:=\{z\in\cp^+;|\Re{z}|<1 \text{ and } \Im{z}>1\}$.
\end{lemma}

\begin{proof}
Let $z,w\in\R$ and let $\S:=S_k$ and $\S_T(x):=T\S(Tx)$ as defined above. Since the map $t\mt G_k(w,z,t)\ft{S_T}(t)$ is not in $C^{\infty}_c(\r)$, we will need an additional trick: Take a function $\p\in C^\infty_c(\r)$, with $\p\geq0$, $\supp(\p)\subset[-1,1]$ and $\ft{\p}(0)=1$. For $0<\ep<1$, let $\p_\ep(x)=\p(x/\ep)/\ep$. Then, for $w,z\in\R$ we define
\begin{equation}
G_{\ep,T}(x)=\left(G_k(w,z,\cd)\ft{\S_T}\right)*\p_\ep(x)
\end{equation}
which belongs to $C^\infty_c(-T(k+1)-1,T(k+1)+1)$. In particular, the Fourier transform $\ft{G_{\ep,T}}(t)$ extends to an entire function. By analytic continuation, for $\zeta\in\U$ we can write
\begin{equation}
\ft{G_{\ep,T}}(\zeta)=\left(\ft{G_k}(w,z,\cd)* \S_T\right)(\zeta)\ft{\p}(\ep \zeta).
\end{equation}
We also note that $\ft{G_{\ep,T}}(\zeta)\ra\ft{G_k}(w,z,\cd)* \S_T(\zeta)$ pointwise for all $\zeta\in\U$ when $\ep\downarrow0$, and we recall the estimate from Lemma \ref{lemma:asymptotic_decay} (ii)
\begin{equation}
\label{eq:unif_estimate_1}
\left|\ft{G_k}(w,z,\cd)* \S_T(\zeta)\right|\ll\frac{1}{(1+(\Re{\zeta})^2)^{k+1}},
\end{equation}
uniformly on $\zeta\in\U$ and $z,w\in\R$. Since $(\mu=\nu+\eta,a)$ is an FS-pair on a strip with $A\subset\U$, we have
\begin{align}
\sum_{|\la|\leq T(k+1)+1} a(\la)G_{\ep,T}(\la)&=\int_{\r} \left(\ft{G_k}(w,z,\cd)* \S_T\right)(t)\ft{\p}(\ep t) d\nu(t)\\&+\sum_{\ga\in A} b(\ga)\left(\ft{G_k}(w,z,\cd)* \S_T\right)(\ga)\ft{\p}(\ep \ga).
\end{align}
Because $a(\cd)$ is locally summable, $G_{\ep,T}\ra G_k(w,z,\cd)\ft{\S_T}$ uniformly on $\r$, $\ft{G_{\ep,T}}(\zeta)\ra\ft{G_k}(w,z,\cd)* \S_T(\zeta)$ pointwise for $\zeta\in\U$ as $\ep\downarrow0$, $|\ft{\p}(\zeta)|\leq c<\infty$ for $\zeta\in\U$ and inequality \eqref{eq:unif_estimate_1}, using the Dominated Convergence Theorem we obtain that
\begin{align}
\label{eq:intermediate_bridge}
\sum_{|\la|\leq T(k+1)} a(\la)G_k(w,z,\la)\ft{\S}\left(\frac{\la}{T}\right)&=\int_{\r} \ft{G_k}(w,z,\cd)*\S_T(t)d\nu(t)\\
&+\sum_{\ga\in A} b(\ga)\left(\ft{G_k}(w,z,\cd)* \S_T\right)(\ga).
\end{align}
holds for any $w,z\in\R$.

Our next step is to take $T\ra\infty$. Observe that the family of maps $\left\{\zeta\in\U\mt\ft{G_k}(w,z,\zeta)\right\}_{z,w\in \R}$ is uniformly continuous and uniformly bounded, uniformly on $z,w\in\R$. Indeed, each function is Lipschitz with constant independent of $z,w$: for $\zeta_1,\zeta_2\in\U$, we have
\begin{align*}
\left|\ft{G_k}(w,z,\zeta_1)-\ft{G_k}(w,z,\zeta_2)\right|&\ll\left|\frac{\ft{G_0}(w,z,\zeta_1)}{(1+\zeta_1^2)^k}-\frac{\ft{G_0}(w,z,\zeta_2)}{(1+\zeta_2^2)^k}\right|\\
&\leq\left|\frac{\ft{G_0}(w,z,\zeta_1)}{(1+\zeta_1^2)^k}-\frac{\ft{G_0}(w,z,\zeta_1)}{(1+\zeta_2^2)^k}\right|+\left|\frac{\ft{G_0}(w,z,\zeta_1)}{(1+\zeta_2^2)^k}-\frac{\ft{G_0}(w,z,\zeta_2)}{(1+\zeta_2^2)^k}\right|\\
&\leq\left|\ft{G_0}(w,z,\zeta_1)\right|\left|\frac{1}{(1+\zeta_1^2)^k}-\frac{1}{(1+\zeta_2^2)^k}\right|\\
&+\left|\ft{G_0}(w,z,\zeta_1)-\ft{G_0}(w,z,\zeta_2)\right|\frac{1}{(1+\zeta_2^2)^k}\\
&\ll |\zeta_1-\zeta_2|,
\end{align*}
because
\begin{align*}
    \left|\ft{G_0}(w,z,\zeta)\right|\leq 2/\pi\, \text{ and } \frac{1}{\left|1+\zeta^2\right|^k}\leq 4^k \text{ uniformly in } \zeta\in\U,z,w\in\R,
\end{align*}
the map $\zeta\in\U\mt\frac{1}{(1+\zeta^2)^k}$ is Lipschitz, and 
\begin{equation}
|\ft{G_0}(w,z,\zeta_1)-\ft{G_0}(w,z,\zeta_2)|=\frac{1}{2\pi|z-\w|}\left|\frac{1}{\zeta_1-z}-\frac{1}{\zeta_1-\w}-\frac{1}{\zeta_2-z}+\frac{1}{\zeta_2-\w}\right|\leq\frac{2}{\pi}\left|\zeta_1-\zeta_2\right|.
\end{equation}
Since $\ft{G_k}(w,z,\cd)$ is uniformly continuous, we deduce that $\ft{G_k}(w,z,\cd)*\S_T(\zeta)\to \ft{G_k}(w,z,\zeta)$ as $T\ra\infty$, for all $\zeta\in\U$. To see that consider the function defined on the horizontal lines $\r+iy$ and reduce to the usual case of $\r$. Combining \eqref{eq:unif_estimate_1} and \eqref{eq:intermediate_bridge}, and using the Dominated Convergence Theorem, we finally deduce that
\begin{align*}
\lim_{T\ra\infty} \sum_{|\la|\leq T(k+1)} a(\la)G_k(w,z,\la)\ft{\S}\left(\frac{\la}{T}\right)&=\frac{1}{2\pi^{k+1}i}\int_{\r} \frac{1}{(t-z)(t-\w)}\cd\frac{d\nu(t)}{(1+t^2)^k}\\
&+\frac{1}{2\pi^{k+1}i}\sum_{\ga\in A} \frac{1}{(\ga-z)(\ga-\w)}\cd\frac{b(\ga)}{(1+\ga^2)^k},
\end{align*}
for each $z,w\in\R$, as desired. 
\end{proof}

\end{document}
